\section{模型检验与敏感性分析}
\subsection{结果口径与可复核性}
三个问题均固定相同的100张计算图和1至5核配置，逐组保存算法方案与赛题评估输出。单核加速比的基准采用题目要求的启发式核内调度结果，而非数学最优值；平均加速比先在每张图上求单核与多核时间之比，再对100张图取算术平均。该值与单核总周期除以多核总周期不同，后者更突出大图的绝对时间贡献，本文同时报告两者，但只将前者作为赛题指定指标。额外搬运量均沿用官方评估输出的定义，不将缓存命中率或代理估计值混入结果。

每个问题的500组方案均经过格式、依赖与官方评估复核。问题一的固定交付证据为\texttt{20260924-A-q1-delivery-r02}；问题二为\texttt{20260924-A-q2-delivery-r07}；问题三为\texttt{20260924-A-q3-delivery-r03}。本文表格和曲线从这些冻结记录生成。对局部优化，保留原方案作为候选，只有官方评价确认改善才替换；这避免由代理排序误差直接造成无条件退步，但不能保证跨场景结果单调，也不能证明全局最优。

\subsection{核数与任务粒度}
图\ref{fig:q1-speed}、\ref{fig:q2-speed}和\ref{fig:q3-speed}表明三种场景的平均加速比总体随核数增加。增核同时提高并行能力，也会增加跨子图边、同步释放点或共享资源竞争，因此单图并非必然单调。问题二记录中有12个相邻核数配置出现较多核更慢的情况；问题三为10个。对这类图，单纯增加核数不是优化动作，应重新联合检查划分、核映射和执行顺序。

子图粒度的作用具有双向性。小子图提供更多可并行任务与调度空隙，但增加跨图数据搬运和任务释放等待；大子图削减边界代价，却可能压缩多核并行度，也可能提高核内活跃数据量并触发换入换出。因此本文生成分量整体、分量内切分及局部边界调整等不同结构的候选，并使用官方完成时间选择，不预设某个算子数量为最优粒度。对大图采用有界候选筛选和局部复用，目的是减少重复核内排程开销；搜索阶段的计时不能代替冷启动、单核基准构建和最终复核在内的端到端计时。

\subsection{共享缓存参数与评价目标}
问题三对典型图进行了L2容量和带宽扫描，而非对100图全部硬件参数做全组合搜索。在案例\texttt{case\_044}的单核检查中，缓存命中可消除部分溢出重读，使完成时间从154407周期降至107814周期；这一大幅收益不应外推至全部图。相反，\texttt{case\_005}五核在已检查的带宽值60、125、250、500下分别得到48200、48202、48259、48253周期，表现为近似持平且略有非单调；\texttt{case\_067}两核的L2方案为32217720周期，高于对应无L2方案的32040130周期。此类负例说明，缓存降低逻辑搬运与缩短关键路径是两个不同命题，L2仲裁、传输重叠和任务释放均须进入评价。

问题三固定配置下，所选L2方案的500组逻辑额外搬运合计比无L2对照多4418288字节；有4组方案时间也高于无L2。本文因此不以减少搬运量或提高命中率单独宣称性能收益。论文报告的结论只对应给定的硬件配置、FIFO规则、官方评价程序和当前100图集合；若换用真实芯片带宽、缓存替换策略或不同图分布，需重新校准并复核。
